\scanpage{10}\mspage{9}
% ink top, lines 1–3
un \uline{demi}-\uline{intervalle} s'il admet un plus petit élément mais pas un plus grand, ou bien l'inverse.

% ink upper middle, NB, lines 1–4
\deleted{\unsure{L'opposé|L'oppo}}\inserted{\uline{NB}} L'opposé d'un tronçon, segment, intervalle, demi-intervalle est itou. L'opposé d'un $1/2$ intervalle fermé à g. est fermé à droite, et inversement.
% ink left, slanting upwards beside NB
% slanted
\margin{Dans ma actuelle terminologie, tronçon \unsure{est|être|était|etre} \unsure{``ens|ens|un|équiv} \uline{biordre}}

% ink middle, Tout, lines 1–3
\deleted{\unsure{Lorsque|Lorsqu.}}\inserted{Tout} tronçon est l'un des trois (segment, $1/2$ intervalle, intervalle) à l'exclusion des deux autres.

% ink lower middle, Si L
Si $\mathcal L$ est un tronçon, on pose $\deleted{\unsure{\partial L}=}$
\[
\begin{aligned}
\unsure{\partial|\partial L}&=\bigl\{\,x\in L\mid
\begin{aligned}&x\text{ est plus petit élément}\\
&\text{ou plus grand élément}\end{aligned}\,\bigr\}
\end{aligned}
\]
\deleted{\illegible{1}}

% ink bottom, Soit and first equality
Soit $\mathcal L$ un tronçon, et
\[
\begin{aligned}
S_0(\mathcal L)=\mathrm{Drap}_2(\mathcal L)&=\{\,\unsure{\varepsilon|\mathcal E|E}\in P_2(\mathcal L)\mid \unsure{\varepsilon|\mathcal E|E}\text{ tot ordonné}\,\}\\
&=\{\,\{x,y\}\mid x<y\,\}
\end{aligned}
\]
% ink bottom, final equality
\[
\text{D'où}\qquad \mathrm{Modf}(\mathcal L)=\mathrm{Drap}(\mathcal L)
=\{\,S\subset P_f(\mathcal L)\mid S\,\substack{\text{tot}\\\text{ord}}\,\}
\]

% pale vertical paper watermark, rotated for reading
\rotatebox{90}{\textsc{\unsure{ra} Strong Extra}}
